% 图契约：四个角部子网的功率递推；图中的端口、几何、箭头和方程对应本题。
% 版面：沿用前三模块的几何图+推导分区；目标印宽17cm，简图至中等密度。
% 布局：左侧真实结构，右侧关键关系；真实电路边不作为架构流程箭头处理。
\documentclass[tikz,border=7pt]{standalone}
\usepackage{fontspec,xeCJK,amsmath,amssymb}
\IfFontExistsTF{Times New Roman}{\setmainfont{Times New Roman}}{\setmainfont{TeX Gyre Termes}}
\IfFontExistsTF{Heiti SC}{\setCJKmainfont[AutoFakeBold=1.5]{Heiti SC}}{\setCJKmainfont[AutoFakeBold=1.5]{FandolHei-Regular}}
\usetikzlibrary{arrows.meta,bending,calc,patterns,positioning,decorations.pathmorphing}
\definecolor{ink}{HTML}{243745}
\definecolor{blue}{HTML}{275DAB}
\definecolor{green}{HTML}{177C63}
\definecolor{red}{HTML}{BD4935}
\definecolor{orange}{HTML}{AA6C16}
\definecolor{purple}{HTML}{715393}
\tikzset{every node/.style={font=\fontsize{9.3}{12.5}\selectfont,text=ink},
arrow/.style={-{Stealth[length=5pt 1.5,width'=0pt 3]},shorten >=2pt,shorten <=1pt,line width=1pt},
arrow short/.style={-{Stealth[length=3pt 1.0,width'=0pt 2.5]},shorten >=1pt,shorten <=0pt,line width=.8pt},
axis/.style={arrow,draw=ink!60,line width=.6pt},
vec/.style={arrow,draw=blue},
wire/.style={draw=ink!75,line width=.8pt},
guide/.style={dashed,draw=ink!45,line width=.5pt},
card/.style={rounded corners=2pt,draw=blue!65,fill=blue!4,align=center,text width=6.5cm,inner xsep=9pt,inner ysep=7pt},
title/.style={font=\fontsize{10.5}{14}\selectfont\bfseries,text width=15cm,align=center},
label/.style={font=\fontsize{8.4}{11}\selectfont,align=center},
note/.style={label,text width=6.1cm}}
\begin{document}
\begin{tikzpicture}[x=1cm,y=1cm]
\path[use as bounding box] (-8,-3.5) rectangle(8,3.65);
\node[title] at (0,3.25) {M5-172d\quad 四个角部子网的功率递推};
\begin{scope}[shift={(-4.05,0)}]\draw[ink](-2,-2)rectangle(-0.6666666666666667,-0.6666666666666667);\draw[ink](0.6666666666666665,-2)rectangle(1.9999999999999998,-0.6666666666666667);\draw[ink](-2,0.6666666666666665)rectangle(-0.6666666666666667,1.9999999999999998);\draw[ink](0.6666666666666665,0.6666666666666665)rectangle(1.9999999999999998,1.9999999999999998);\draw[ink,thick](-0.6666666666666667,-2)--(0.6666666666666665,-2)(-0.6666666666666667,2)--(0.6666666666666665,2)(-2,-0.6666666666666667)--(-2,0.6666666666666665)(2,-0.6666666666666667)--(2,0.6666666666666665);\fill[red] (-2,-2)circle(2pt);\node[left=4pt] at(-2,-2){$A$};
\fill[red] (2,-2)circle(2pt);\node[right=4pt] at(2,-2){$B$};
\fill[red] (2,2)circle(2pt);\node[right=4pt] at(2,2){$C$};
\fill[red] (-2,2)circle(2pt);\node[left=4pt] at(-2,2){$D$};
\draw[orange,very thick](-.6667,-2)--(.6667,-2);\node[orange,below]at(0,-2.05){$l/3$ 中段};\node[above]at(0,2.3){四角子网按 $1/3$ 缩放};\end{scope}\node[card,draw=blue!65,fill=blue!4]at(3.9,1.8){$P(l)=4[P(l/3)+R_{AC}(l/3)I_c^2+(lr/3)I_c^2]$};
\node[card,draw=green!65,fill=green!4]at(3.9,0){$\mathcal E(l/3)=\mathcal E(l)/9$};
\node[card,draw=orange!65,fill=orange!4]at(3.9,-1.8){$P(l)=343k^2l^3/(1104r)$};

\end{tikzpicture}\end{document}